\section{A connected minimal sufficient dependency audit}
\label{sec:dependency-audit-v30}

\status{The stipulated correctness of the checked Lean proof is accepted.
This audit concerns only completion of its human reconstruction.
It connects the final theorem to the actual common-family and host-count
obligations, separates ordinary RH from unrestricted Dirichlet GRH, and
classifies every retained continuation record.  Minimality here is
relative to the explicitly selected source route after removing proved
implications, duplicate producers and stronger optional interfaces.  It
is not a proof that no mathematically shorter proof exists, nor a
transitive enumeration of every Lean declaration.}

\subsection{How to read the audit}

The classification of an argument and its completion status are separate.
A deep published theorem may be a completed application of class A; an
elementary assertion may be an open class-B task if its original objects
have not yet been identified.  In particular, a proved implication
$P\Rightarrow Q$ is not a proof of $Q$ while $P$ is still a producer
obligation.

\begin{center}
\begin{tabular}{p{.08\linewidth}p{.83\linewidth}}
A&Direct application of a published theorem, with the exact consumed
hypotheses checked.  The certificate records its theorem, normalization,
source objects and parameter range.\\
B&Elementary or moderately difficult argument that can be reconstructed
independently, including finite algebra, numerical inequalities and
precise assembly on the original objects.\\
C&Genuinely substantial mathematical reconstruction of the actual
arithmetic, automorphic, geometric or spectral producer.\\
D&A possible gap in our understanding of the source: an unresolved
definition, object identity, normalization, theorem match or quantifier
transport.  This is not an assertion of a gap in the formal proof.\\
\end{tabular}
\end{center}

The new machine-readable map has explicit logical edges for the chosen
source implications.  Work-sharing and possible-reuse edges are separately
marked: identifying a common local dictionary is useful to several
producers, but the dictionary alone is not asserted to prove any of them.
The complete records and definitions remain attached to the frontier.

This repairs a navigation defect in version 0.29.  Following only its
logical \texttt{requires}/alternative edges from the two final roots
reached 28 of its 151 nodes; the count/cohomology branch was attached
through explanatory \texttt{related\_nodes}, not a logical host edge.
The earlier map was a sufficient-interface guide, but could not be used
as a connected completion audit.  The earlier map and every proof are
preserved as historical records.

\subsection{The exact conditional route to the final theorem}

Fix a primitive target $\chi$ first.  The selected route is:
\[
 \begin{gathered}
 \text{original count suppliers}
       \Longrightarrow \mathrm{VWinAt}(\chi,401)
       \mathrel{\mathop{\Longrightarrow}^{\rm HX4}}
                 \mathrm{HostExW}(3/2,(\mathrm{GChi}\chi).\mathrm{req}),\\
 \text{same-target host}+
       \text{one compatible design and its closed links}
       \Longrightarrow \mathrm{FamilyClassConnected},\\
 \mathrm{FamilyClassConnected}
       \mathrel{\mathop{\Longrightarrow}^{\rm ZW4,WB8,AN3}}
                L(s,\chi)\ne0\quad(\Re s>1/2),
 \end{gathered}                                                    \tag{DA30.1}
\]
with the principal pole treated as in the exported statement.
Every arrow is quantified over the original source models and preserves
the target; the host window alone does not construct the compatible
geometric family.  Applying this route at $\chi=1$ gives ordinary RH.
Applying it for every primitive target, and then GH6's imprimitive-factor
and functional-equation argument, gives Dirichlet GRH.

The common-family interface is exactly the argument list of
\texttt{topDesWin}.  It asks for a compatible design and the following
nine closed predicates on the same class, together with its class
inclusions and \texttt{HypV}:
\[
\begin{gathered}
 \mathrm{LinkN},\ \mathrm{LinkEta},\ \mathrm{LinkIRR},\
 \mathrm{SmallLevelRankEq},\
 (\mathrm{TuSize}\Rightarrow\mathrm{DGamSize}\Rightarrow\mathrm{LinkDIVt}),\\
 (\mathrm{ChosenE}\Rightarrow\mathrm{NuFullF}\wedge\mathrm{MuRank}),\
 \mathrm{LinkG},\
 (\mathrm{TuSize}\Rightarrow\mathrm{LinkS}),\\
 (\mathrm{TuSize}\Rightarrow\mathrm{DGamSize}\Rightarrow\mathrm{LinkGLOBAL}).
\end{gathered}                                                    \tag{DA30.2}
\]
Thus there are exactly ten first-level producer interfaces: the design
and these nine closed predicates.  In the selected \texttt{shapeWin}
caller, $(C_{\rm abs},C_0,e_D,e_L)=(2,3,2,0)$; design existence is on
$\mathrm{WithTw}(\mathrm{WithSm}\,\mathcal C')$, the nine predicates
are closed on $\mathrm{Des}(\mathcal C,3)$, and the additional guards
are $\mathcal C'\subseteq\mathcal C\cap\mathcal E$ and
$\mathcal E\Rightarrow\mathrm{HypV}$.  Here ``closed'' means that for
every sufficiently small positive $\epsilon$ one threshold is chosen
before all families and height parameters in that class; the predicate
then holds whenever the family is large at that threshold.  Design
existence must work for every such transported threshold.

The finite threshold maximum is taken before design existence selects
one family.  It would be invalid to supply each predicate on an
independently chosen family.  The target class gives both the required
$\mathrm{EAtH}$ and the same cohomological $\mathrm{HypV}$.
DA1--DA2 prove the assembly; MT7 now discharges its previously retained
metric premise.  The thresholds, coefficient fields, word sets, lattices,
local flags, maps and metrics keep their literal bindings.

\subsection{The smaller counting cut for ordinary RH}

For an unrestricted target, KA4 and the existing YX/HV/TW/OB/VW chain
give the sufficient cohomological frontier
\[
 \mathrm{CorGK}_{K_\chi},\quad \mathrm{PMExAlt},\quad
 \mathrm{CountIII}_\chi,\quad \mathrm{CountV}_\chi,                \tag{DA30.3}
\]
relative to their recorded classical and prime-pair inputs.
Here $K_\chi=K_*(p_K(\chi))$ is chosen once for the fixed target;
$\mathrm{PropGK}_{K_\chi}$ is a sufficient stronger input for its
first entry.  The quantifiers over later hosts, coefficient DVRs,
closures, finite prime sets and layer functions remain intact.

KA5 proves a genuine target specialization.  Every original host
character $\psi$ satisfies
$\operatorname{ord}(\psi)\mid2^{s+2}\operatorname{ord}(\chi)$.
Case III requires $\operatorname{ord}(\psi)$ to be odd and
$\psi\ne1$.  If $\operatorname{ord}(\chi)$ is a power of two,
these conditions contradict one another.  KA6 checks that each consumer
through \texttt{VWinAt} uses only that same fixed target.
Consequently (DA30.3) reduces to
\[
 \boxed{\quad
 \mathrm{CorGK}_{K_\chi},\quad\mathrm{PMExAlt},\quad
 \mathrm{CountV}_\chi\qquad
        (\operatorname{ord}(\chi)\text{ a power of two}).\quad}    \tag{DA30.4}
\]
In particular, the separate Case-III seed and promodular construction
do not lie on this selected route to ordinary RH.  They remain on the
unrestricted Dirichlet route.  No universal \texttt{CountIII} theorem
has been inferred from a single-target empty case.

For $\chi=1$, $K_\chi\cong\mathbb Q(\sqrt{-7})$ and the exceptional
set is $\{3,7\}$; KA gives the reduced-form proof of its class number.
Count V still concerns unbounded dyadic conductors at those two primes.
KA's concrete examples show that Cases I, II and V remain nonempty,
including characteristic three.  A theorem requiring $\ell>3$ therefore
cannot automatically close the entire principal support branch.

\subsection{What this version actually proves}

\begin{center}
\begin{longtable}{p{.12\linewidth}p{.53\linewidth}p{.24\linewidth}}
Claims&Completed argument&Boundary retained\\\hline
MT1--MT7&The original ideal relation, exact slope change, all floor-error
support costs and one uniform constant give literal \texttt{LinkMu}.&
Actual \texttt{LinkDIV}, fullness and positive rank; the family itself.\\
HX1--HX4&All old host inputs are installed on the original field and
prime pair; the cohomological rescaling preserves $401\to400$.&
$\mathrm{VWinAt}(\chi,401)$; exact classical prime-pair leaves.\\
KA1--KA4&All Case-I \texttt{GD} fields, two-prime valuation aggregation,
finite kernel-orbit mass and original coefficient return give
$\mathrm{PropGK}\Rightarrow\mathrm{CorGK}\Rightarrow\mathrm{CountI}$.&
The actual uniform congruence integer/one-character producer.\\
KA5--KA6&Case-III vacuity at a power-of-two-order target propagates through
every height and window consumer.&Three remaining cohomological inputs
in (DA30.4); the separate global family.\\
ZW1--ZW4&The full weighted-slope and numerical error argument maps literal
\texttt{ChainTied} into the exact $1/2+(e_D+e_L+7)\epsilon$ window.&
Production of the arithmetic ledger upstream.\\
AQ1--AQ5&An arbitrarily late positive-part integral bound is enough for
zero exclusion; integer values suffice for pointwise windows after
elementary interpolation.&No new arithmetic supplier of either weaker
contract is claimed.\\
\end{longtable}
\end{center}

The metric proof uses the actual identity $N_u=I_uM_u$ already stored
in \texttt{FamP}, with
$I_u=\prod_w\mathfrak p_w^{-\lfloor e_ws_{u,w}\rfloor}$.
The comparison metric is unchanged by definition, so
\[
 \mu(N_u)-\mu(M_u)
       -2\sum_w\ell_we_ws_{u,w}
      =2\sum_w\ell_w(\lfloor e_ws_{u,w}\rfloor-e_ws_{u,w}).
\]
After every actual exceptional prime family is included, MT proves the
uniform constant $2(\log4+5+10|C_\Delta|)$ on the original $X''$ scale.
This removes an entire independent interface; it does not supply the
upstream divisibility or spectral theorem.

ZW similarly removes a numerical obligation without inventing an
arithmetic shortcut.  If $E$ is the original error bracket and
$R=1+|K|+|c_{HH}|+|e_K|$, its proof gives
\[
 \mathrm{ErrX5at}\le200R^3E\,x^{(e_D+e_L+2)\epsilon}
                  \le E\,x^{(e_D+e_L+3)\epsilon}
\]
after one uniform threshold.  WB then gives the additional four epsilon
units.  The geometric slope inequality is a supplied ledger field;
the finite numerical deduction from that field is now fully written.

\subsection{The remaining global obligations}

The table classifies the still-unsupplied original interfaces; many
conditional sublemmas are already complete.  Shared dependencies listed
in the right column are work-sharing information, not automatically
sufficient hypotheses for the entire left-hand contract.
The dictionary and row-system entries refine the ten interfaces in
(DA30.2); they are not extra independent first-level inputs.

\begin{center}
\begin{longtable}{p{.22\linewidth}p{.08\linewidth}p{.59\linewidth}}
Original obligation&Class&Exact remaining work\\\hline
Compatible design&C&One actual Core/layout, coefficient embeddings,
$\mathrm{Over/RecOK/LedOK}$, rows and lattices, with all requested
theta/local data.  Initial zero theta or empty toric data are not the
required final installation.\\
Local/theta/operator dictionary&D&Identify the actual maps, local models,
nonzero theta lifts and critical toric ratios used by the same
$\mathrm{NTh/NChoice/ChoiceSide/LocAll}$ family.  Their remaining
nonvanishing, convergence, separation and comparison arguments are C.\\
Actual rows and spectrum&C&Original $\mathrm{St/ColPin/R3.Read}$,
rational integral realization, retained rows, and original operator
actions.  Canonical representation irreducibility alone does not
construct these.\\
$\mathrm{LinkN}$&C&Actual local Galois denominator table and its exact
type/theta/rank bindings.  LN proves the full conditional family chain,
including the restoration cost $d_j\log|G|$.\\
$\mathrm{LinkEta}$&C&Compatible twist coverage at every actual local
place; nonempty places alone are insufficient.\\
$\mathrm{LinkIRR}$&C&One original row system with injective label count,
all joins and row count, yielding the precise retained defect.\\
Small-level rank&D&Locate and match the actual projection-surjectivity
supplier $z_L$ and original $\mu_0$; the subsequent finite-dimensional
argument is B.\\
$\mathrm{LinkDIVt}$&C&Original rational period, specified isometries,
heights, ideal shifts, support and word-height variation, all on the
same target and operators.\\
Fullness/rank and global slope&D/B/C&Match the common real-place second
layer, mass ratio and genus-volume normalization.  The same source
\texttt{link\_gl} feeds both interfaces; count this arithmetic
matching once.  Object identification is D, finite installation is B,
and a substantive comparison not already supplied remains C.\\
$\mathrm{LinkG}$&C&Simultaneous original lattice lengths, orbit coverage,
mean-kappa bound and uniform gain; previously proved local statements
must use the one actual global instance.\\
$\mathrm{LinkS}$&C&Normalized finite-word, Poisson, image/norm, kernel,
mass and Gram identities with the exact $\Gamma$, $D_*$, $D_{\rm all}$
and exponential error terms.\\
\end{longtable}
\end{center}

The original class-inclusion and target-class definitions are B
identifications, recorded separately from the C construction of their
objects.  These definitions and the complete argument lists occur in
the pinned source-anchor inventory.  No large package is assigned A
merely because its name contains ``literature'' or a familiar theorem.

\subsection{The remaining counting obligations}

\begin{center}
\begin{longtable}{p{.20\linewidth}p{.08\linewidth}p{.61\linewidth}}
Branch&Class&Original unresolved producer and smaller substeps\\\hline
Count I&C&\texttt{OneChar} requires actual unramified-coefficient
automorphy and the Kato input.  The rank-three/rank-two congruence theorem
must choose one nonzero integer before every later coefficient prime.
Its final orbit/Frobenius/valuation installation into \texttt{PropGK}
is B.  Exact Kato and polarized-adjoint literature matching is D until
the original representations and hypotheses are certified.  The
already matched QX adjoint application does not itself certify this
different rank-three use.\\
Count II&C&On each of QX's two actual Links, one and the same ideal must
satisfy all six \texttt{HidaPropsAlt} fields and \texttt{P1core}, above
one common coefficient threshold.  QX and YX already give
$\mathrm{PMExAlt}\Rightarrow\mathrm{CountII}$.
Matching a candidate published support theorem is D.\\
Count III&C&For unrestricted targets: actual \texttt{SeedIVAlt},
tensor-factor tower, \texttt{ChainRankFour}, residual descent and seeded
chain reduction, then the full original Case-III promodular witness.
Finite tower/descent portions are B.  This separate branch is absent
from (DA30.4), while the common support input remains.\\
Count V&C/D&Original rank-three support/automorphy \texttt{ThmS6\_s},
host rigidity and uniform host catalogue.  The original frame/length
installation, \texttt{HostVanE}, and \texttt{HostXpp} are B after the
already supplied cohomology and CorX inputs.  Rank-four statements
cannot be imported without a checked D bridge.\\
\end{longtable}
\end{center}

The detailed count inventory records the exact order of all existential
and universal choices.  Some of its edges are explicitly only proposed
reuse paths: for example, a separate proof of the rank-three catalogue
does not require a rank-four-to-rank-three bridge.  Likewise separate
existence of a Hida ideal and a support ideal does not give their
required common witness.

\subsection{Published inputs and stronger optional routes}

Seven applications were freshly checked against primary sources in this
round: the quadratic analytic class-number normalization; Bombieri's
exceptional-factor density theorem; Heath--Brown's hybrid density
specialization; Thorne's exact polarized adjoint vanishing; Caraiani's
conjugate-self-dual local compatibility and temperedness; Hida's
rank-two fixed-central control application; and complete-local
connectedness after the dimension bounds.  HX1 adds the precisely
specialized least-prime theorem.  Their applications are A in these
scopes, not broader arithmetic existence theorems.

The literature inventory also retains 101 entries from 22 earlier
classical-input registries.  These were inventoried, not all independently
rechecked in this round.  Their original written hypothesis certificates
remain the authority for the exact scope.  An empty convenience-summary
field in the new inventory does not erase the detailed original record.
Where an old certificate is reused for a different rank, polarization,
local condition or coefficient-prime range, that new match is D.

In particular, the rank-two Hida input used for the Case-II seed does not
supply the rank-four six-field Hecke ideal.  Allen--Newton--Thorne's
residually reducible automorphy theorem is a candidate support source,
but its residual-Schur, constituent and local hypotheses and its
$\ell>3$ range have not been matched to both actual Links.  It is not
an automatic A closure of \texttt{P1core}.

The backlog keeps optional stronger routes separate: universal
Harish--Chandra interfaces beyond the selected large-gap route,
internal proofs of already matched named classical theorems, all-rank
counting packages, unchanged-level versions, and the stronger one-level
symmetric predicate.  Their absence does not become an additional
requirement of the selected sufficient route.  Every one of the 245
inherited continuation records has an explicit disposition, including
compound records whose closed and still-open subscopes differ.

The complete disposition contains 245 records in 107 scope groups.
Their primary classes are A: 3, B: 10, C: 118 and D: 114; independently,
215 records remain applicable to the combined RH/Dirichlet-GRH project,
25 concern optional stronger routes, and 5 are closed in their stated
scope.  These are bookkeeping counts of inherited records, with overlap
and compound scopes, not counts of independent mathematical obligations
or measures of proof completion.  In particular, ``applicable'' does not
mean required separately for ordinary RH: the Case-III specialization
above still applies.  The three class-A records in this continuation
inventory are not the same inventory as the seven freshly rechecked
published applications plus HX1.

\subsection{A weaker terminal contract worth keeping available}

AQ1--AQ4 show that pointwise square-root control is stronger than the
analytic endpoint needs.  For each fixed target and each $\theta$ in
a sequence decreasing to $1/2$, it is enough to fix one $\delta>0$
such that for every $\eta>0$ and every $Y_0$ there is
$Y\ge\max(2,Y_0)$ with
\[
 \int_Y^{Y^{1+\delta}}\max(H_\chi(x),0)x^{-\theta-1}\,dx
                        \le Y^\eta.                              \tag{DA30.5}
\]
The positive part is essential; a signed integral could hide a positive
peak through cancellation.  The family and $\delta$ precede $\eta$
and the choice of window.  AQ's gamma localization and nearest-pole
argument prove the full implication, including possible real poles.

AQ5 also shows that integer upper bounds on a fixed-stretch window
are enough: on $[n,n+1]$ the affine interpolation $L_n$ satisfies
$0\le H_\chi-L_n\le1/(8n)$, because
$H_\chi''(x)=-1/x$ there.  Slightly shrinking the stretch keeps the
endpoints inside the integer window.

Neither observation currently supplies an arithmetic producer.  The
audit therefore records (DA30.5) as an alternative research contract,
not as a completed shortcut around the common-family obligations.
The older ``direct gap'' contract is treated the same way: bundling
the unresolved arithmetic into one predicate does not reduce its work.

\subsection{The next branch by contribution to completion}

The next default investigation should be the \emph{actual shared
local/theta/operator installation}, with the shared real-place
second-layer/mass normalization investigated in parallel when possible.
These inputs are used by design existence, denominator comparisons,
period/height control, local gains and the spectral theorem.  Resolving
their exact source dictionary is a concrete D task and determines which
large C proofs really remain.  Its output must be a literal common
instance and an exact downstream consumer; another conditional theorem
assuming that instance would not retire the producer.

Among bounded follow-ups, the actual level-projection surjectivity
$z_L$ is a promising target once its supplier is located.  The original
Count-V frame/length and vanishing installations, and the
\texttt{OneChar} plus rank-three-number to \texttt{PropGK} join, are
other finite B chains.  Their value is to expose or remove whole
consumer interfaces, not to increase the number of lemma tags.

The rank-four ordinary support construction remains a substantial
high-value task on both RH and GRH routes, but it is no longer the
automatic next choice solely because version 0.29 ended there.
The separate Case-III seed/promodular route is a priority for completing
unrestricted Dirichlet GRH after taking account of its absence from
the selected ordinary-RH route.

The full human-readable proof is still unfinished.  This version gives a
connected, classified route to that goal, completes five useful groups
of arguments, and makes the remaining original-object obligations
reviewable without changing the assumption about the checked formal
proof.  Source and text checks in this round do not constitute a new
Lean replay.
